% Topic document: official format, 11 pt and 1-inch margins.
\documentclass[11pt]{article}
\usepackage[utf8]{inputenc}
\usepackage[T1]{fontenc}
\usepackage[margin=1in]{geometry}
\usepackage{amsmath,amssymb}
\usepackage{xcolor}
\usepackage[round]{natbib}
\usepackage[colorlinks=true,allcolors=blue!50!black]{hyperref}
\title{Reservas bancarias, colateral y asignaci\'on de capital}
\author{Alejandro Ventura \\ \small\url{https://github.com/alejandroventuraventurameza-tech/ai-project}}
\date{Propuesta de tema: 7 de octubre de 2026}
\begin{document}
\maketitle

\section{La pregunta y el track}
\textbf{Track B: modelo de tesis.} \textquestiondown Bajo qu\'e condiciones una
expansi\'on de reservas reduce el costo del cr\'edito y mejora la asignaci\'on de
un capital fijo entre empresas heterog\'eneas en productividad, patrimonio y
colateral? El alcance acordado \textbf{A} es un n\'ucleo est\'atico con un banco
y dos empresas. El alcance \textbf{B}, objetivo de la presentaci\'on final, ampl\'ia
la transmisi\'on y los reg\'imenes. Estas letras designan alcances del proyecto;
la elecci\'on del curso sigue siendo Track B. La manufactura formal peruana
motiva la pregunta, pero esta propuesta no identifica efectos causales.

\section{El baseline y el modelo que falta}
El antecedente principal es \citet{gonzalez2026}, versi\'on de enero de 2026:
ya conecta pol\'itica monetaria, heterogeneidad y mala asignaci\'on. La
restricci\'on es $q_t k_t\leq\gamma q_t a_t$ (ec.~4, p.~7); el problema est\'atico
maximiza $m_t f_t-w_t l_t-R_t k_t$ (ec.~6, p.~8), con $m_t$ precio real
del insumo y $R_t$ costo del capital. Esa notaci\'on no es nuestro $m,R$.
Sus proposiciones 1 y 2 (pp.~19 y 20--21; pruebas B.7.1--B.7.2) separan cambios
en la distribuci\'on patrimonial y en el umbral de productividad. Nuestro modelo
no rederiva esas proposiciones: usa dos empresas con rendimientos decrecientes,
costos de liquidez bancaria y colateral terminal heterog\'eneo, frente al
continuo con rendimientos constantes del antecedente.

\citet{bianchi2022} aporta reservas, dep\'ositos, pr\'estamos y liquidaci\'on
interbancaria con fricciones. El archivo local es el manuscrito de marzo de
2021, no la maquetaci\'on publicada: su secci\'on 2 describe las dos etapas y las
ecuaciones 2--3 (p.~6) el balance y la restricci\'on de capital bancario.
El costo cuadr\'atico siguiente es una \textbf{simplificaci\'on propia}; no es
una ecuaci\'on ni una proposici\'on de Bianchi y Bigio.

\paragraph{Banco e instrumento de cantidad.}
Con reservas $m$ dadas por la autoridad, patrimonio bancario $E\geq0$, cr\'edito
$B$ y dep\'ositos $D$, el banco toma los retornos brutos $t,\rho,r_m$ y resuelve
\[
 \max_{B,D\geq0}\ tB+r_m m-\rho D-C(B,D,m),\qquad B+m=D+E,
\]
\[
 C=cB+\frac{\kappa}{2}(\delta D-m)_+^2,\quad
 c\geq0,\quad\kappa>0,\quad0<\delta<1,\quad (v)_+=\max\{v,0\}.
\]
Todas las posiciones se expresan en unidades del bien inicial; los pagos y el
costo, en el bien terminal. $\delta D$ resume necesidades de liquidaci\'on,
no un mercado OTC expl\'icito. En un \mbox{\emph{interior}}, con liquidez escasa,
\[
 t=\rho+c+\kappa\delta[\delta(B+m-E)-m].
\]
\newpage
El experimento A var\'ia $m$ manteniendo $\rho$ fijo. Un hogar pasivo ofrece
dep\'ositos y capital; no se resuelve su bienestar ni un equilibrio monetario
din\'amico completo. El precio de reservas $r_m$ afecta beneficios, pero, con
$m$ fijo, no aparece en esta condici\'on marginal.

\paragraph{Empresas, unidades y timing.}
Al inicio cada empresa $i\in\{H,L\}$ posee fondos $n_i>0$, usa $x_i$ y toma
deuda principal $b_i\geq0$ para alquilar $k_i\geq0$ a precio $R>0$. Produce al
final $y_i=A_i k_i^\alpha$, con $A_H>A_L>0$, $0<\alpha<1$, y paga $tb_i$.
El l\'imite $h_i>0$ es riqueza terminal pignorable ex\'ogena; no es el capital
adquirido. Los fondos no usados rinden $\rho>0$. Omitiendo la dotaci\'on
terminal constante, el problema es
\begin{equation}\label{eq:firm}
 \max_{k_i,x_i,b_i\geq0}\ A_i k_i^\alpha-tb_i+\rho(n_i-x_i),
 \quad Rk_i=x_i+b_i,\quad x_i\leq n_i,\quad tb_i\leq h_i.
\end{equation}
La desigualdad $t>\rho$ elimina el arbitraje de endeudarse para ahorrar:
$x_i=\min\{Rk_i,n_i\}$ y $b_i=(Rk_i-n_i)_+$. Por tanto,
\[
 \max_{0\leq k_i\leq\bar k_i}\ A_i k_i^\alpha-\rho Rk_i
 -(t-\rho)(Rk_i-n_i)_+ +\rho n_i,
 \qquad\bar k_i=\frac{n_i+h_i/t}{R}.
\]
Este objetivo es estrictamente c\'oncavo. Definiendo
$M_i=\alpha A_i k_i^{\alpha-1}$, sus condiciones marginales son
\[
\begin{array}{ll}
 M_i=\rho R,& Rk_i<n_i \quad\text{(autofinanciamiento interior)},\\
 \rho R\leq M_i\leq tR,& Rk_i=n_i \quad\text{(quiebre, sin deuda)},\\
 M_i=tR,& n_i<Rk_i<n_i+h_i/t\quad\text{(deuda interior)},\\
 M_i=tR(1+\mu_i),\quad\mu_i\geq0,
 &Rk_i=n_i+h_i/t\quad\text{(colateral activo)}.
\end{array}
\]
Aqu\'i $\mu_i$ multiplica $h_i-tb_i$ en el problema original. En el interior
con deuda, la condici\'on de $x_i$ hace usar todos los fondos propios; con deuda
limitada, $M_i>tR$ significa $\mu_i>0$. Los casos de igualdad delimitan cambios
de r\'egimen y no admiten trasladar derivadas de un interior sin comprobaci\'on.

\paragraph{Cierre y medida de mala asignaci\'on.}
Se fija la oferta f\'isica $K>0$, $k_H+k_L=K$, y se determina $R$ por vaciado
del mercado. El cr\'edito bancario debe satisfacer $B=b_H+b_L$ y su balance,
con $D\geq0$. Una reducci\'on de $t$ a $R$ fijo no prueba una mejora de
asignaci\'on: hay que resolver tambi\'en $R$. El benchmark sin fricciones es
\[
 k_H^*=K\frac{A_H^{1/(1-\alpha)}}{A_H^{1/(1-\alpha)}+A_L^{1/(1-\alpha)}},
 \quad Y^*=A_H(k_H^*)^\alpha+A_L(K-k_H^*)^\alpha.
\]
Se usar\'an $\mathcal L=1-Y/Y^*$ y
$G=|\log M_H-\log M_L|$, donde $Y=A_Hk_H^\alpha+A_Lk_L^\alpha$.
La igualdad $M_H=M_L$ caracteriza la asignaci\'on eficiente interior. Como
$K$ permanece fijo, una subida de $Y$ implica una subida de productividad
agregada a insumos dados; no es el contrafactual emp\'irico de Hsieh--Klenow.

\newpage
\section{El resultado esperado}
Los siguientes son \textbf{resultados candidatos con derivaci\'on preliminar},
no una afirmaci\'on universal sobre pol\'itica monetaria. Se mantienen fijos
$K,A_i,n_i,h_i,\rho,E,c,\kappa,\delta$. Toda comparaci\'on se restringe a un
intervalo donde se conservan las desigualdades estrictas del r\'egimen.

\paragraph{P1: H limitada y L se autofinancia.}
Sup\'ongase $b_H=h_H/t$, $M_H>tR$ y $Rk_L<n_L$, de modo que
$M_L=\rho R$. Si $s=1/(1-\alpha)$, el equilibrio real satisface
\[
 \frac{n_H+h_H/t}{R}+\left(\frac{\alpha A_L}{\rho R}\right)^s=K.
\]
La derivaci\'on impl\'icita entrega
\[
 \frac{dR}{dt}=-\frac{h_H}{t^2(k_H+sk_L)}<0,\qquad
 \frac{dk_H}{dt}=-\frac{sk_Lh_H}{t^2R(k_H+sk_L)}<0.
\]
Como $M_H>tR>\rho R=M_L$, H est\'a subcapitalizada respecto de $k_H^*$.
Una ca\'ida local de $t$ reasigna capital desde L hacia H, eleva $Y$ y reduce
$\mathcal L$ y $G$. El resultado exige que H siga limitada y L siga en
autofinanciamiento interior. No cubre el quiebre $Rk_L=n_L$.

\paragraph{P2: ambas empresas limitadas.}
Sup\'ongase $b_i=h_i/t$ y $M_i>tR$ para ambas. El vaciado del mercado implica
\[
 R=\frac{n_H+n_L+(h_H+h_L)/t}{K},\qquad
 k_H=K\frac{tn_H+h_H}{t(n_H+n_L)+h_H+h_L}.
\]
Entonces
\[
 \frac{dk_H}{dt}
 =\frac{K(n_Hh_L-n_Lh_H)}{[t(n_H+n_L)+h_H+h_L]^2}.
\]
Una ca\'ida de $t$ aumenta $k_H$ \textbf{si y solo si}
$h_H/n_H>h_L/n_L$; la igualdad produce neutralidad en las participaciones.
Para que esa reasignaci\'on mejore eficiencia se necesita adem\'as
$k_H<k_H^*$ y que no se cruce $k_H^*$ en la comparaci\'on finita. Si H ya
est\'a sobrecapitalizada, darle m\'as capital puede empeorar la asignaci\'on.
El signo del efecto sobre producto se verifica mediante
\[
 \frac{dY}{dt}=(M_H-M_L)\frac{dk_H}{dt}.
\]

\paragraph{De reservas a costo del cr\'edito.}
En P1, $B(t)=h_H/t$; en P2, $B(t)=(h_H+h_L)/t$. En ambos $B'(t)<0$.
Al imponer simult\'aneamente el equilibrio bancario en liquidez escasa,
\[
 \frac{dt}{dm}=\frac{-\kappa\delta(1-\delta)}{1-\kappa\delta^2B'(t)}<0.
\]
La derivada incluye la respuesta del cr\'edito. Con reservas abundantes,
$t=\rho+c$ y m\'as reservas son neutrales. El caso $t=\rho$ y un encaje
activo requieren otra caracterizaci\'on.

\newpage
\section{Por qu\'e no est\'a ya resuelto}
La conexi\'on entre dinero y mala asignaci\'on \textbf{ya existe}. Se usa el
manuscrito local de enero de 2026 (173 p\'aginas PDF): ecuaciones 4 y 6,
proposiciones 1--2 y pruebas B.7.1--B.7.2. Las versiones de 2023 y BdE 2021
quedan como contexto hist\'orico. La versi\'on elegida tiene tres autores:
Gonz\'alez, Nu\~no y Thaler; reconoce a Albrizio y retira el bloque
emp\'irico anterior. La web del autor informa una revisi\'on
solicitada en JPE: aqu\'i se cita como manuscrito, no como art\'iculo publicado.
Tambi\'en se localiz\'o el trabajo relacionado \emph{A tale of two margins:
Monetary policy and capital misallocation} (European Economic Review, 2026);
su existencia refuerza la necesidad de delimitar el aporte, no certifica novedad.

El aporte por evaluar es una descomposici\'on de la transmisi\'on desde liquidez
bancaria hacia capital relativo, con $tb_i\leq h_i$ y dos reg\'imenes. La
condici\'on $h_H/n_H>h_L/n_L$ depende de este cierre. No se replica la
din\'amica distributiva ni el problema Ramsey de 2026. La b\'usqueda de citantes
sigue parcial; no se afirma haber probado originalidad.

\section{Plan y riesgos}
\textbf{Tema:} cerrar A, verificar las cuatro regiones, P1--P2 y el equilibrio
real y bancario, y entregar propuesta y diapositivas compiladas.
\textbf{Para el 30 de octubre:} desarrollar B como extensi\'on controlada:
separar instrumento de precio, reservas y encaje; estudiar transiciones de
r\'egimen y, si es tractable, acumulaci\'on de patrimonio. No se promete un
DSGE completo, dolarizaci\'on ni una pol\'itica \`optima.

\textbf{Lean, despu\'es de fijar el art\'iculo:} integrar el PDF del manuscrito,
registrar su commit y ejecutar la habilidad de formalizaci\'on de
AppliedModelingLib sobre ese PDF inmutable, con \texttt{gpt-5.6-sol} y esfuerzo
\texttt{xhigh}. Objetivos: identidad de P2, signo de su cambio finito, y P1
con todas las hip\'otesis econ\'omicas declaradas. Ejecutar el chequeo por
art\'iculo \texttt{--fast}, conservar su salida y la carpeta generada completa
en \texttt{lean/} respetando su
\texttt{.gitignore}, y mapear cada resultado a declaraci\'on, archivo y estado.
Sin \texttt{sorry}, sin predicados vac\'ios ni hip\'otesis que asuman la conclusi\'on.
Repetir si cambia una proposici\'on. Lean sigue pendiente.

\textbf{Riesgo principal:} la sensibilidad del resultado al colateral
ex\'ogeno, al capital fijo y a cambios de r\'egimen. La comprobaci\'on num\'erica
usar\'a sorteos y contraejemplos; no sustituye la demostraci\'on.
EEA 2023 y ENAHO 2025 permanecer\'an como motivaci\'on descriptiva pendiente de
auditor\'ia. No se usar\'an sus proxies o ganancias preliminares de TFP como
identificaci\'on causal ni como calibraci\'on validada.
\bibliographystyle{plainnat}
\bibliography{../paper/references}
\end{document}
